% =====================================================================================
\section{Conclusion}
\label{sec:conclusion}
% =====================================================================================

This report examined whether the heat that spreads between the waveguides of a thermo-optic phased array, normally
suppressed, can be used to steer it. In simulation it can, provided the structure makes the steady temperature
profile take the shape of the required phase.

Heat spreading acts as a spatial low-pass filter, and the ratio of spreading length to pitch determines what a few
heaters can do. On a standard stack ($\Lth/d\approx1.7$) a few heaters cannot steer. On a membrane undercut beneath
the array ($\Lth/d\approx10$) three or four heaters steer at about \qty{4.9}{\mWdeg}, but most of the edge heater's
power flows into the nearby wall. A slot past that heater removes this path, so the heat crosses the array and the
temperature becomes a linear ramp. Three heaters then steer \num{33} emitters across \ang{22.8} at \qty{1.38}{\mWdeg}
in air, and one heater does so at \qty{0.53}{\mWdeg} in vacuum, within \qty{5}{\percent} of a conduction limit that is
independent of the number of emitters. Uniform heating of a membrane between two walls gives a parabolic profile and
hence a one-wire tunable lens at about \qty{1}{\milli\watt\per\rad}.

The three capabilities of individual control are lost to a quantified extent. Without phase wrapping the power grows
with angle, the hottest waveguide is about \qty{62}{\kelvin} above ambient at \ang{10}, and a \ang{10} step takes
\qtyrange{2}{9}{\milli\second} with a plain drive, or \qtyrange{0.16}{1.4}{\milli\second} when each heater's
drive is shaped in time, provided the thermal time scale is known to a few per cent. Only two to four phase patterns are usable. Independent fabrication errors cannot be
corrected, so a one-time static trim is required, after which calibration reduces to a linear fit and a gain curve.

Compared with published devices, the design needs \numrange{10}{30} times less power than measured few-wire arrays and
\numrange{30}{100} times less than measured few-heater lenses, and is correspondingly slower. It is not the
lowest-power phased array: \num{33} of the best folded phase shifters would hold \ang{10} with about a quarter of the
power, but would require \num{33} wires and a fan-in.

These results rest on two-dimensional finite-element models that pass their convergence, conservation and consistency
checks, and on an array model that reproduces its analytic limits. They have not been tested in three dimensions or on
a fabricated device. Within these limits, thermal crosstalk can be used as a resource for the smooth phase profiles
that steering and focusing require, at a quantified cost in speed and control.

\section*{Data and code availability}
\addcontentsline{toc}{section}{Data and code availability}
The COMSOL exports, the Python code, the result tables and the figure scripts are available at
\url{https://github.com/shanmehta2210/Diffusion_in_Waveguides}.
